\begin{table}[htbp]
\centering
\caption{Common organisational structures for product development teams}
\label{tab:session2-organisational-structures}
\renewcommand{\arraystretch}{1.2}
\begin{tabularx}{\textwidth}{|>{\raggedright\arraybackslash}p{0.18\textwidth}|>{\raggedright\arraybackslash}X|>{\raggedright\arraybackslash}p{0.26\textwidth}|}
\hline
\textbf{Structure} & \textbf{Main characteristic} & \textbf{Typical use in APID terms} \\
\hline
Functional & Work stays mainly within specialised roles or disciplines. & Suitable when the project is simple and roles are stable. \\
\hline
Project & The team is organised mainly around the project itself. & Suitable when the group needs fast decisions and strong coordination. \\
\hline
Lightweight & A project lead coordinates, but most authority remains with functional roles. & Suitable when the work is shared but not highly integrated. \\
\hline
Heavyweight & The project lead has stronger authority across functions and priorities. & Suitable when integration, speed, and cross-functional decisions matter greatly. \\
\hline
\end{tabularx}
\end{table}